\SourceLeaf{286}%
\noindent where $F$ denotes the area of the whole figure and $S$ that of the entire surface,
\[
= (2\chi - \sigma) \, . \, 360^0 - \frac{F}{S} \,.\,720^0.
\]
This rule includes the one given above, since the unbounded region in the plane has character 0, making $\frac{F}{S} = 0$. Above, the signs were counted where the intersection points were encountered for the first time, and there they are opposite to those used here.

\begin{center}
[7.]
\end{center}

\begin{center}\textbf{ANOTHER WAY TO TREAT THESE QUESTIONS.}\end{center}

Let $S$ be the whole curve and $s$ an indeterminate portion. Draw a tangent at the starting point and denote by $(0, s)$ the angle made by the straight line from the starting point to the endpoint of the portion $s$. Clearly, $(0, s)$ will vary continuously from $s = 0$ as long as the curve has not returned to the starting point. We assume that this does not happen until $s = S$, and the last value before $s$ reaches this value will be
\[
(0, s) = \pm 180^0,
\]
according as the angles are taken to increase toward the inner or the outer side of the line.

With $s$ held constant, draw a straight line from each point of the line, regarded as the endpoint of the smaller portion $s'$, to the endpoint of the portion $s$, and call its inclination to the first tangent $(s', s)$. As long as the endpoint of $s'$ does not coincide with that of $s$, varying $s'$ continuously will make $(s', s)$ vary continuously, provided that the portion $s$ does not intersect itself, until $s' = s$. At this limiting value itself, $(s', s)$ is strictly speaking indeterminate, but as long as $s' < s$ it is determinate however close the values may be. The closer $s'$ comes to $s$, the more closely this value approaches the inclination of the tangent with respect to the tangent at the starting point. Now let $s = S - \omega$, and let $s'$ vary from $0$ to $S - \omega - \omega'$. Clearly the change in $s'$ will thus come all the closer to $180^0$ the smaller $\omega, \omega'$ are, so that $(S - \omega - \omega', S - \omega)$ comes
